% ============================================================================
% APPENDIX JT: LIT-TIROLE - Jean Tirole Research Integration
% ============================================================================
% Category: LIT (Literature Integration)
% Version: 1.0
% Last Updated: 2026-01-17
% Dependencies: K (LIT-FEHR), GB (LIT-BECKER), IOO (DOMAIN-IO)
% Language: english
% ============================================================================

% ----------------------------------------------------------------------------
% HEADER BLOCK
% ----------------------------------------------------------------------------
% Code: JT
% Category: LIT (Literature Integration)
% Title: Jean Tirole Research Integration
% Author: EBF Framework
% Status: Complete
% Priority: High
% Evidence-Base: 60 papers (see data/paper-sources.yaml)
% ----------------------------------------------------------------------------

\chapter{LIT-TIROLE: Jean Tirole Research Integration}
\label{app:lit-tirole}


\begin{tcolorbox}[colback=blue!5!white, colframe=blue!75!black,
    title=Quick Reference: Appendix JT]

\textbf{Appendix:} JT (LIT)

\textbf{Core Question:} What are the key findings in this domain?

\textbf{Key Concepts:}
\begin{itemize}[nosep]
\item Framework integration
\item Parameter validation
\item Cross-references
\end{itemize}

\textbf{Cross-References:} See related appendices below
\end{tcolorbox}

% ============================================================================
% CROSS-REFERENCE MAP
% ============================================================================
\begin{tcolorbox}[colback=blue!5!white,colframe=blue!75!black,title=Cross-Reference Map]
\textbf{Outgoing References:}
\begin{itemize}
    \item CORE-WHO (AAA): Multi-stakeholder welfare in corporate governance
    \item CORE-WHAT (C): Non-monetary utility dimensions (intrinsic motivation, image)
    \item CORE-HOW (B): Platform complementarities and network effects
    \item CORE-WHEN (V): Context-dependent regulation design
    \item CORE-WHERE (BBB): Market structure parameters
    \item CORE-AWARE (AU): Information asymmetries in regulation
    \item CORE-READY (AV): Willingness under incentive structures
    \item CORE-STAGE (AW): Market evolution and platform lifecycle
    \item LIT-FEHR (K): Social preferences and crowding out
    \item LIT-BENABOU (RB): Behavioral IO collaboration (motivation, narratives)
    \item LIT-FALK (FAL): Moral reasoning collaboration (narratives paper)
    \item LIT-BECKER (GB): Rational choice benchmark
    \item LIT-THALER (M): Behavioral insights in market design
    \item DOMAIN-IO (IOO): Industrial organization applications
    \item METHOD-LLMMC (AN): Parameter estimation for market power
\end{itemize}

\textbf{Incoming References:}
\begin{itemize}
    \item DOMAIN-IO (IOO): Market power and competition theory
    \item DOMAIN-FIN (AF): Corporate finance and governance
    \item FORMAL-GAME (D): Game-theoretic foundations
    \item PREDICT-MARKET (AS): Market structure predictions
\end{itemize}
\end{tcolorbox}

% ============================================================================
% CHAPTER LINKAGE
% ============================================================================
\begin{tcolorbox}[colback=green!5!white,colframe=green!75!black,title=Chapter Linkage]
\textbf{Primary Chapter:} Chapter 9 (Market Design and Institutions)

\textbf{Supporting Chapters:}
\begin{itemize}
    \item Chapter 3: Incentive theory foundations
    \item Chapter 5: Information economics
    \item Chapter 10: Corporate governance and stakeholder welfare
    \item Chapter 14: Platform economics applications
\end{itemize}
\end{tcolorbox}

% ============================================================================
% ABSTRACT AND QUICK REFERENCE
% ============================================================================
\begin{abstract}
This appendix integrates Jean Tirole's foundational contributions to economics into the EBF framework. Tirole (Nobel Prize 2014) revolutionized our understanding of market power, regulation, platform economics, and the interaction between intrinsic and extrinsic motivation. His work with Roland Bénabou on behavioral industrial organization provides crucial bridges between traditional IO and behavioral economics. The EBF framework incorporates Tirole's insights through six axioms (JT-1 to JT-6) that formalize two-sided market dynamics, regulatory design under asymmetric information, motivation crowding, and corporate governance. Key contributions include the theory of two-sided markets ($\pi = (p_1 - c_1)D_1 + (p_2 - c_2)D_2$ with $D_i(p_1, p_2)$), the Laffont-Tirole regulatory framework, and the Bénabou-Tirole model of intrinsic motivation crowding.
\end{abstract}

\begin{tcolorbox}[colback=gray!5!white,colframe=gray!75!black,title=Quick Reference]
\textbf{Core Contributions:}
\begin{itemize}
    \item Industrial organization: Market power, entry barriers, price discrimination
    \item Two-sided markets: Platform economics, network effects, interchange fees
    \item Regulation theory: Adverse selection, moral hazard, regulatory capture
    \item Behavioral IO: Intrinsic motivation, identity, self-confidence
    \item Corporate governance: Stakeholder theory, control rights, delegation
    \item Contract theory: Incomplete contracts, renegotiation, authority
\end{itemize}

\textbf{Key Parameters:}
\begin{itemize}
    \item Platform pricing: $p_1^* + p_2^* \neq 2 \cdot MC$ (asymmetric pricing)
    \item Motivation crowding: $\partial e / \partial b < 0$ when $b$ signals distrust
    \item Regulatory information rent: $\Delta\theta \cdot q(\bar{\theta})$ under adverse selection
    \item Network effect strength: $\partial D_i / \partial n_j > 0$ (cross-side effects)
\end{itemize}

\textbf{Evidence Tier:} Mixed (Tier 2-5)
\begin{itemize}
    \item Platform economics: Tier 2-3 (quasi-experiments, case studies)
    \item Motivation crowding: Tier 1-2 (lab experiments, field studies)
    \item Regulatory theory: Tier 3-5 (case studies, theoretical)
    \item Corporate governance: Tier 2-3 (event studies, cross-country)
\end{itemize}
\end{tcolorbox}

% ============================================================================
% FUNDAMENTAL QUESTION
% ============================================================================
% -----------------------------------------------------------------------------
% APPENDIX SCOPE BOX
% -----------------------------------------------------------------------------

\begin{tcolorbox}[
    colback=orange!5!white,
    colframe=orange!75!black,
    title={\textbf{Appendix Scope: Ziel / In-Scope / Out-of-Scope / Constraints}},
    fonttitle=\bfseries
]

\textbf{Ziel (Objective):}
\begin{itemize}[nosep]
    \item How should we design markets, regulations, and incentive structures when agents have private information, intrinsic motivations, and operate in networked environments with multiple stakeholders?
\end{itemize}

\textbf{In-Scope:}
\begin{itemize}[nosep]
    \item Fundamental Question
    \item Key Results and Evidence
    \item EBF Integration
    \item Worked Example: Platform Pricing Strategy
    \item Critical Assessment
\end{itemize}

\textbf{Out-of-Scope (delegiert an andere Appendices):}
\begin{itemize}[nosep]
    \item REF-GLOSSARY $\rightarrow$ Appendix G
    \item REF-GLOSSARY $\rightarrow$ Appendix G
    \item DOMAIN-IO $\rightarrow$ Appendix IOO
\end{itemize}

\textbf{Constraints (Anwendungsgrenzen):}
\begin{itemize}[nosep]
    \item [TODO: Define when this appendix does NOT apply]
    \item [TODO: Define required prerequisites]
    \item [TODO: Define known limitations]
\end{itemize}

\textbf{Lieferobjekte (Deliverables):}
\begin{itemize}[nosep]
    \item 4 Table(s)
    \item 14 Equation(s)
    \item 6 Axiom(s)
    \item Worked Example(s)
\end{itemize}

\end{tcolorbox}


\section{Fundamental Question}
\label{sec:jt-fundamental}

\begin{tcolorbox}[colback=red!5!white,colframe=red!75!black,title=Fundamental Question]
\textbf{How should we design markets, regulations, and incentive structures when agents have private information, intrinsic motivations, and operate in networked environments with multiple stakeholders?}
\end{tcolorbox}

Jean Tirole's research program addresses the fundamental challenge of designing economic institutions---markets, firms, regulatory bodies---when the standard assumptions of perfect competition and complete information fail. His work provides:

\begin{enumerate}
    \item Tools for analyzing market power and designing competition policy
    \item Frameworks for regulating industries with private information
    \item Understanding of platform dynamics in the digital economy
    \item Integration of behavioral insights into industrial organization
    \item Multi-stakeholder perspectives on corporate governance
\end{enumerate}

For the EBF framework, Tirole's contributions are essential because they:
\begin{itemize}
    \item Bridge traditional IO with behavioral economics (Bénabou-Tirole)
    \item Provide tools for analyzing context-dependent market outcomes
    \item Formalize how incentive structures interact with intrinsic motivation
    \item Offer frameworks for multi-level welfare analysis (platforms, regulators, users)
\end{itemize}

% ============================================================================
% THEORETICAL FOUNDATIONS
% ============================================================================
\section{Theoretical Foundations}
\label{sec:jt-theory}

\subsection{Two-Sided Markets and Platform Economics}
\label{subsec:jt-platforms}

Tirole's work with Rochet \citep{tirole_rochet_2003_platform, tirole_rochet_2006_twosided} revolutionized our understanding of platforms:

\begin{equation}
\pi = (p_1 - c_1) \cdot D_1(p_1, p_2, n_2) + (p_2 - c_2) \cdot D_2(p_1, p_2, n_1)
\label{eq:jt-platform}
\end{equation}

where $D_i$ depends on prices to both sides and the number of participants on the other side ($n_j$).

\begin{axiom}[Two-Sided Market Pricing]
\label{ax:jt-1}
\textbf{JT-1:} In two-sided markets, optimal pricing does not equate price to marginal cost on each side. Instead, platforms subsidize the side with stronger cross-side externalities:
\begin{equation}
p_i^* = c_i - \frac{\partial D_j / \partial n_i}{\partial D_i / \partial p_i} \cdot (p_j - c_j)
\end{equation}
The ``subsidy'' side is determined by relative elasticities and network effects, not absolute costs.
\end{axiom}

\textbf{EBF Integration:} Axiom JT-1 connects to CORE-HOW (B) through complementarity---user value depends on other-side participation. It connects to CORE-WHO (AAA) through multi-stakeholder welfare: platform, side 1 users, side 2 users, and society.

\textbf{Key Insight:} The price level ($p_1 + p_2$) affects market coverage; the price structure ($p_1/p_2$) affects which side is subsidized. Competition policy must consider both.

\subsection{Regulation Under Asymmetric Information}
\label{subsec:jt-regulation}

The Laffont-Tirole framework \citep{tirole_laffont_1993_regulation} addresses regulation when firms have private information:

\begin{equation}
\max_{q, t} \int [S(q(\theta)) - (1+\lambda)(c(\theta)q(\theta) + t(\theta))] dF(\theta)
\label{eq:jt-regulation}
\end{equation}

subject to incentive compatibility (IC) and participation (IR) constraints.

\begin{axiom}[Information Rent vs. Allocative Efficiency]
\label{ax:jt-2}
\textbf{JT-2:} Optimal regulation trades off allocative efficiency against information rents:
\begin{equation}
\text{Distortion} \propto \frac{\lambda}{1+\lambda} \cdot \frac{F(\theta)}{f(\theta)}
\end{equation}
where $\lambda$ is the shadow cost of public funds. High-cost types face binding IC constraints; low-cost types earn information rents. Bunching occurs when the virtual cost function is non-monotonic.
\end{axiom}

\textbf{EBF Integration:} Axiom JT-2 connects to CORE-AWARE (AU) through information asymmetries and to CORE-WHEN (V) through context-dependent regulatory design. The key insight: one-size-fits-all regulation is suboptimal when firms differ in unobservable characteristics.

\subsection{Regulatory Capture}
\label{subsec:jt-capture}

Tirole's work on regulatory capture \citep{PAP-tirole_laffont_1988_politics_regulation} formalizes how interest groups influence regulation:

\begin{equation}
\Pr(\text{capture}) = g(I - k)
\label{eq:jt-capture}
\end{equation}

where $I$ is the regulated firm's stake and $k$ is the ``bribe'' cost.

\begin{axiom}[Regulatory Capture Dynamics]
\label{ax:jt-3}
\textbf{JT-3:} Regulatory capture is more likely when:
\begin{enumerate}
    \item Stakes are concentrated (high $I$ for few firms)
    \item Monitoring costs are high (low $k$ for bribes)
    \item Regulators have career concerns in the industry
    \item Information asymmetries favor incumbents
\end{enumerate}
Institutional design (separation of powers, transparency, term limits) can reduce capture vulnerability.
\end{axiom}

\textbf{EBF Integration:} Axiom JT-3 connects to CORE-WHO (AAA) through the multi-level welfare problem: regulator welfare vs. public welfare vs. firm welfare.

\subsection{Intrinsic and Extrinsic Motivation}
\label{subsec:jt-motivation}

The Bénabou-Tirole framework \citep{tirole_benabou_2003_intrinsic_motivation, tirole_benabou_2006_incentives_prosocial} integrates behavioral insights into incentive theory:

\begin{equation}
e^* = \arg\max_e \left[ v_a \cdot a(e) + v_y \cdot y(e) + b \cdot e - C(e) \right]
\label{eq:jt-motivation}
\end{equation}

where $v_a$ is intrinsic motivation, $v_y$ is reputational concern, and $b$ is the explicit bonus.

\begin{axiom}[Motivation Crowding]
\label{ax:jt-4}
\textbf{JT-4:} Extrinsic incentives can crowd out intrinsic motivation when they signal the principal's beliefs about the agent:
\begin{equation}
\frac{\partial e^*}{\partial b} = \underbrace{\frac{\partial e^*}{\partial b}\bigg|_{v_a}}_{\text{direct } > 0} + \underbrace{\frac{\partial e^*}{\partial v_a} \cdot \frac{\partial \hat{v}_a}{\partial b}}_{\text{indirect (crowding)}}
\end{equation}
Crowding out occurs when $\partial \hat{v}_a / \partial b < 0$---i.e., high bonuses signal that the task is unpleasant or that the principal distrusts the agent.
\end{axiom}

\textbf{EBF Integration:} Axiom JT-4 is central to EBF's understanding of incentive design. It connects to CORE-WHAT (C) by showing utility dimensions beyond monetary payoffs, and to CORE-HOW (B) through the complementarity between intrinsic and extrinsic motivation (which can be positive or negative depending on context).

\textbf{Key Insight:} The same bonus can increase or decrease effort depending on what it signals about the task or the relationship.

\subsection{Identity and Motivated Beliefs}
\label{subsec:jt-identity}

The Bénabou-Tirole work on identity \citep{tirole_benabou_2011_identity, tirole_benabou_2016_mindful_economics} shows how beliefs serve as assets:

\begin{equation}
V(\hat{\theta}) = \max_e \left[ U(e, \theta) + \beta \cdot \mathbb{E}_{\hat{\theta}}[V(\hat{\theta}')] \right]
\label{eq:jt-identity}
\end{equation}

where $\hat{\theta}$ represents self-beliefs about ability or morality.

\begin{axiom}[Beliefs as Assets]
\label{ax:jt-5}
\textbf{JT-5:} Agents may distort beliefs when beliefs provide anticipatory utility or motivation benefits:
\begin{equation}
\hat{\theta}^* = \arg\max_{\hat{\theta}} \left[ \underbrace{U(\theta, \hat{\theta})}_{\text{decision utility}} + \underbrace{\mu(\hat{\theta})}_{\text{anticipatory}} - \underbrace{\kappa(\hat{\theta}, \theta)}_{\text{reality costs}} \right]
\end{equation}
This explains overconfidence, denial, and the persistence of motivated beliefs in the face of contradicting evidence.
\end{axiom}

\textbf{EBF Integration:} Axiom JT-5 connects to CORE-AWARE (AU) through the gap between objective information and subjective awareness. It explains why information alone often fails to change behavior---beliefs are maintained because they serve psychological functions.

\subsection{Corporate Governance and Stakeholders}
\label{subsec:jt-governance}

Tirole's work on corporate governance \citep{PAP-tirole_2001_corporate_governance} extends beyond shareholder value:

\begin{equation}
W = \sum_{i \in \text{stakeholders}} \omega_i \cdot U_i(a)
\label{eq:jt-governance}
\end{equation}

\begin{axiom}[Multi-Stakeholder Governance]
\label{ax:jt-6}
\textbf{JT-6:} Optimal corporate governance balances stakeholder interests based on:
\begin{enumerate}
    \item \textbf{Contractibility:} Can the stakeholder's interest be protected contractually?
    \item \textbf{Specificity:} Is the stakeholder's investment relationship-specific?
    \item \textbf{Measurability:} Can the stakeholder's contribution be measured?
    \item \textbf{Voice/Exit:} Does the stakeholder have voice or exit options?
\end{enumerate}
Shareholders get control rights not because they are most important, but because their interests are hardest to protect contractually.
\end{axiom}

\textbf{EBF Integration:} Axiom JT-6 connects directly to CORE-WHO (AAA) through the welfare hierarchy---firm-level welfare aggregates employee, shareholder, customer, and community welfare.

% ============================================================================
% KEY RESULTS AND EVIDENCE
% ============================================================================
\section{Key Results and Evidence}
\label{sec:jt-results}

\subsection{Platform Economics Evidence}

\begin{table}[htbp]
\centering
\caption{Two-Sided Market Evidence}
\label{tab:jt-platform-evidence}
\begin{tabular}{llcc}
\toprule
\textbf{Platform} & \textbf{Finding} & \textbf{JT-1 Support} & \textbf{Evidence Tier} \\
\midrule
Credit cards & Merchant fees $\gg$ consumer fees & Yes (subsidy side) & Tier 2 \\
Gaming consoles & Console sold below cost & Yes (hardware subsidy) & Tier 3 \\
Search engines & Free to users, ads to businesses & Yes (user subsidy) & Tier 3 \\
Ride-sharing & Driver and rider subsidies & Yes (growth phase) & Tier 2 \\
Dating apps & Women often free, men pay & Yes (scarce side) & Tier 3 \\
\bottomrule
\end{tabular}
\end{table}

\subsection{Motivation Crowding Evidence}

\begin{table}[htbp]
\centering
\caption{Intrinsic Motivation Crowding: Empirical Evidence}
\label{tab:jt-crowding-evidence}
\begin{tabular}{llcc}
\toprule
\textbf{Study} & \textbf{Finding} & \textbf{JT-4 Support} & \textbf{Evidence Tier} \\
\midrule
Gneezy \& Rustichini (2000) & Small fines increased late pickups & Yes (crowding out) & Tier 1 \\
Frey \& Oberholzer-Gee (1997) & Payment reduced NIMBY acceptance & Yes (civic duty) & Tier 2 \\
Ariely et al. (2009) & High bonuses hurt performance & Yes (pressure) & Tier 1 \\
Mellström \& Johannesson (2008) & Payment reduced blood donation & Yes (image) & Tier 1 \\
Bénabou \& Tirole (2006) & Image concerns explain prosociality & Yes (theoretical) & Tier 2 \\
\bottomrule
\end{tabular}
\end{table}

\subsection{Regulatory Evidence}

\begin{table}[htbp]
\centering
\caption{Regulation and Capture: Case Evidence}
\label{tab:jt-regulation-evidence}
\begin{tabular}{llcc}
\toprule
\textbf{Case} & \textbf{Finding} & \textbf{JT-2/3 Support} & \textbf{Evidence Tier} \\
\midrule
US telecom deregulation & Incumbent advantages persisted & Yes (information rents) & Tier 3 \\
Financial regulation 2008 & Regulatory forbearance & Yes (capture) & Tier 3 \\
UK utility regulation & Price caps with ratcheting & Yes (dynamic IC) & Tier 3 \\
EU antitrust (Big Tech) & Platform-specific remedies & Yes (two-sided) & Tier 3 \\
\bottomrule
\end{tabular}
\end{table}

% ============================================================================
% EBF INTEGRATION
% ============================================================================
\section{EBF Integration}
\label{sec:jt-integration}

\subsection{CORE Framework Mapping}

\begin{table}[htbp]
\centering
\caption{Tirole's Contributions Mapped to 10C Framework}
\label{tab:jt-9c-mapping}
\begin{tabular}{lll}
\toprule
\textbf{CORE} & \textbf{Tirole Contribution} & \textbf{Application} \\
\midrule
WHO (AAA) & Multi-stakeholder governance & Platform welfare, corporate governance \\
WHAT (C) & Non-monetary utility (image, intrinsic) & Motivation crowding \\
HOW (B) & Network complementarities & Two-sided markets \\
WHEN (V) & Context-dependent regulation & Optimal mechanism design \\
WHERE (BBB) & Market structure parameters & Markup estimation, network effects \\
AWARE (AU) & Information asymmetries & Adverse selection in regulation \\
READY (AV) & Incentive compatibility & Participation constraints \\
STAGE (AW) & Platform lifecycle & Growth vs. monetization phases \\
\bottomrule
\end{tabular}
\end{table}

\subsection{Behavioral Industrial Organization}

Tirole's collaboration with Bénabou creates ``Behavioral IO''---integrating behavioral economics into market analysis:

\begin{tcolorbox}[colback=green!5!white,colframe=green!75!black,title=Behavioral IO: The Bénabou-Tirole Program]
\textbf{Traditional IO assumes:}
\begin{itemize}
    \item Agents respond only to material incentives
    \item Preferences are stable and exogenous
    \item Information affects beliefs but not preferences
\end{itemize}

\textbf{Behavioral IO adds:}
\begin{itemize}
    \item Intrinsic motivation as utility component
    \item Image concerns and social signaling
    \item Motivated beliefs and identity maintenance
    \item Incentive-belief interactions (crowding)
\end{itemize}

\textbf{Policy implications:}
\begin{itemize}
    \item Incentive design must consider signaling effects
    \item Information campaigns may backfire if they threaten identity
    \item Market design should preserve intrinsic motivation where valuable
\end{itemize}
\end{tcolorbox}

% ============================================================================
% WORKED EXAMPLE
% ============================================================================
\section{Worked Example: Platform Pricing Strategy}
\label{sec:jt-example}

\begin{tcolorbox}[colback=yellow!5!white,colframe=yellow!75!black,title=Worked Example: Ride-Sharing Platform Pricing]
\textbf{Scenario:} A ride-sharing platform (like Uber/Lyft) must set prices for riders ($p_R$) and commission rates for drivers ($p_D$). The platform faces:
\begin{itemize}
    \item Rider demand: $D_R = 100 - 2p_R + 0.5n_D$ (more drivers → more riders)
    \item Driver supply: $n_D = 50 - p_D + 0.3D_R$ (more riders → more drivers)
    \item Marginal cost per ride: $c = 5$
\end{itemize}

\textbf{Standard Pricing (ignore network effects):}
\begin{equation}
p_R = p_D = c + \text{markup} = 5 + 3 = 8
\end{equation}

\textbf{Tirole Two-Sided Pricing (JT-1):}
\begin{enumerate}
    \item Solve for equilibrium participation:
    \begin{align}
    D_R^* &= 100 - 2p_R + 0.5(50 - p_D + 0.3D_R^*) \\
    n_D^* &= 50 - p_D + 0.3D_R^*
    \end{align}

    \item Recognize cross-side externalities:
    \begin{equation}
    \frac{\partial D_R}{\partial n_D} = 0.5, \quad \frac{\partial n_D}{\partial D_R} = 0.3
    \end{equation}

    \item Optimal pricing subsidizes the driver side (stronger network contribution):
    \begin{equation}
    p_R^* = 10, \quad p_D^* = 4 \quad (\text{below cost!})
    \end{equation}
\end{enumerate}

\textbf{EBF Extension (Behavioral IO):}
Consider driver intrinsic motivation (flexibility, autonomy). If high commission signals distrust:
\begin{equation}
\text{Driver effort} = e(p_D, \hat{v}_a(p_D))
\end{equation}
where $\partial \hat{v}_a / \partial p_D < 0$ (high fees signal platform doesn't value drivers).

\textbf{Optimal strategy:} Lower commission + transparency about platform costs to preserve driver intrinsic motivation and service quality.

\textbf{Policy Implication:} Regulating only rider prices (price caps) may backfire by forcing platforms to extract more from drivers, reducing service quality.
\end{tcolorbox}

% ============================================================================
% CRITICAL ASSESSMENT
% ============================================================================
\section{Critical Assessment}
\label{sec:jt-critical}

\subsection{Strengths of Tirole's Framework}

\begin{enumerate}
    \item \textbf{Analytical rigor:} Game-theoretic foundations enable precise predictions
    \item \textbf{Policy relevance:} Direct applications to antitrust and regulation
    \item \textbf{Behavioral integration:} Bénabou-Tirole work bridges IO and behavioral economics
    \item \textbf{Multi-stakeholder view:} Goes beyond shareholder value maximization
    \item \textbf{Dynamic considerations:} Addresses platform evolution and regulatory adaptation
\end{enumerate}

\subsection{Limitations and Extensions}

\begin{enumerate}
    \item \textbf{Equilibrium focus:} Less attention to out-of-equilibrium dynamics
    \item \textbf{Rationality assumptions:} Behavioral IO is still largely rational-agent based
    \item \textbf{Empirical calibration:} Many parameters are hard to estimate
    \item \textbf{Distributional effects:} Limited attention to inequality within platforms
    \item \textbf{Data and privacy:} Digital platform issues not fully addressed in early work
\end{enumerate}

% ============================================================================
% SUMMARY
% ============================================================================
\section{Summary}
\label{sec:jt-summary}

\begin{tcolorbox}[colback=blue!5!white,colframe=blue!75!black,title=Key Takeaways]
\begin{enumerate}
    \item \textbf{JT-1 (Two-Sided Markets):} Platform pricing subsidizes the side with stronger cross-side externalities
    \item \textbf{JT-2 (Regulation):} Optimal regulation trades off efficiency against information rents
    \item \textbf{JT-3 (Capture):} Concentrated interests and information asymmetries enable capture
    \item \textbf{JT-4 (Crowding):} Extrinsic incentives can crowd out intrinsic motivation via signaling
    \item \textbf{JT-5 (Beliefs):} Agents maintain motivated beliefs that serve psychological functions
    \item \textbf{JT-6 (Governance):} Multi-stakeholder governance allocates control based on contractibility
\end{enumerate}

\textbf{Integration Principle:} Tirole's framework extends rational choice to settings with market power, asymmetric information, and networked interactions. The Bénabou-Tirole behavioral extensions show how incentive design must account for non-material motivations.
\end{tcolorbox}

% ============================================================================
% GLOSSARY
% ============================================================================
\section{Glossary}
\label{sec:jt-glossary}

Key terms defined in Appendix G (REF-GLOSSARY):

\begin{itemize}
    \item \textbf{Two-sided market:} Market where platform serves two distinct user groups with cross-side externalities
    \item \textbf{Network effect:} Value to each user increases with number of other users
    \item \textbf{Cross-side externality:} Benefit to one side from participation of the other side
    \item \textbf{Information rent:} Profit earned from private information in mechanism design
    \item \textbf{Regulatory capture:} When regulatory agency serves regulated industry rather than public
    \item \textbf{Motivation crowding:} When extrinsic incentives reduce intrinsic motivation
    \item \textbf{Motivated beliefs:} Beliefs distorted to serve psychological needs (anticipatory utility, self-image)
\end{itemize}

% ============================================================================
% CRITICAL FOUNDATIONS
% ============================================================================
\section{Critical Foundations}
\label{sec:jt-foundations}

\begin{enumerate}
    \item Tirole, J. (1988). \textit{The Theory of Industrial Organization}. MIT Press.
    \item Laffont, J.-J., \& Tirole, J. (1993). \textit{A Theory of Incentives in Procurement and Regulation}. MIT Press.
    \item Rochet, J.-C., \& Tirole, J. (2003). Platform Competition in Two-Sided Markets. \textit{JEEA}, 1(4), 990--1029.
    \item Rochet, J.-C., \& Tirole, J. (2006). Two-Sided Markets: A Progress Report. \textit{RAND Journal}, 37(3), 645--667.
    \item Bénabou, R., \& Tirole, J. (2003). Intrinsic and Extrinsic Motivation. \textit{Review of Economic Studies}, 70(3), 489--520.
    \item Bénabou, R., \& Tirole, J. (2006). Incentives and Prosocial Behavior. \textit{AER}, 96(5), 1652--1678.
    \item Bénabou, R., \& Tirole, J. (2011). Identity, Morals, and Taboos: Beliefs as Assets. \textit{QJE}, 126(2), 805--855.
    \item Tirole, J. (2001). Corporate Governance. \textit{Econometrica}, 69(1), 1--35.
    \item Tirole, J. (2015). Market Failures and Public Policy (Nobel Lecture). \textit{AER}, 105(6), 1665--1682.
    \item Fudenberg, D., \& Tirole, J. (1991). \textit{Game Theory}. MIT Press.
\end{enumerate}

% ============================================================================
% OPEN ISSUES
% ============================================================================
\section{Open Issues}
\label{sec:jt-open}

\begin{enumerate}
    \item \textbf{Platform regulation:} How to regulate digital platforms with data network effects?
    \item \textbf{Algorithmic pricing:} Does AI-enabled pricing facilitate tacit collusion?
    \item \textbf{Gig economy:} How to apply stakeholder governance when workers are ``independent''?
    \item \textbf{Attention markets:} How to value attention as the scarce resource in ad-supported platforms?
    \item \textbf{Cross-platform competition:} How to analyze competition between non-comparable platforms?
\end{enumerate}

% ============================================================================
% REFERENCES SECTION
% ============================================================================
\section{References}
\label{sec:jt-references}

\nocite{bcm_master}

This appendix integrates 60 papers from Jean Tirole's research corpus. Full references are maintained in the master bibliography (\texttt{bibliography/bcm\_master.bib}).

\textbf{Core Tirole Works (60 papers):}
\begin{itemize}
    \item \textbf{Industrial Organization:} Tirole (1988), Hart \& Tirole (1990), Fudenberg \& Tirole (1983, 1984, 1985)
    \item \textbf{Regulation:} Laffont \& Tirole (1991, 1993, 1994, 1996, 1999, 2000)
    \item \textbf{Platforms:} Rochet \& Tirole (2003, 2006, 2008), Caillaud \& Tirole (2003), Armstrong \& Tirole (2007)
    \item \textbf{Behavioral IO:} Bénabou \& Tirole (2002, 2003, 2004, 2006, 2009, 2010, 2011, 2016); Bénabou, Falk \& Tirole (2018: Narratives)
    \item \textbf{Corporate Governance:} Tirole (2001), Holmström \& Tirole (1991, 1997, 1998, 2000)
    \item \textbf{Game Theory:} Fudenberg \& Tirole (1991), Maskin \& Tirole (1999)
    \item \textbf{Innovation:} Aghion \& Tirole (1997), Lerner \& Tirole (2002, 2004, 2005)
    \item \textbf{Banking \& Liquidity:} Dewatripont \& Tirole (1994), Holmström \& Tirole (1998), Tirole (2011), Farhi \& Tirole (2012, 2021)
    \item \textbf{Antitrust:} Rey \& Tirole (1986, 2007)
    \item \textbf{Climate:} Tirole (2010), Weitzman \& Tirole (2019)
    \item \textbf{Digital Economy:} Tirole (2021)
    \item \textbf{Macroeconomics:} Tirole (1982, 1985)
\end{itemize}

\textbf{Key Empirical Applications:}
\begin{itemize}
    \item Gneezy \& Rustichini (2000): Motivation crowding in childcare
    \item Ariely et al. (2009): Performance under high incentives
    \item Card \& Krueger (1994): Minimum wage and employment
\end{itemize}

See Appendix G (REF-GLOSSARY) for term definitions and Appendix IOO (DOMAIN-IO) for industrial organization applications.

\end{document}
